
\section{Conclusions and Future Directions}\label{sec discuss}
This paper sheds light on under-explored phenomena in pruning practices for neural network model compression. On a theoretical level, we prove an accurate distributional characterization (DC) for the solution of overparameterized least-squares for linear models with correlated Gaussian features. Our DC allows to precisely characterize the pruning performance of popular pruning methods, such as magnitude pruning. Importantly, our DC combined with a linear Gaussian equivalence, leads to precise analytic formulas for the pruning performance of nonlinear random feature models. On the experimental side, we provide a thorough study of overparameterization and pruning with experiments on linear models, random features and neural nets with growing complexity. Our experiments reveal striking phenomena such as a novel double descent behavior for model pruning and the power of overparameterization. They also shed light on common practices such as retraining after pruning. 
% \cts{Finish}

Going forward, there are several exciting directions to pursue. First, it would be insightful to study whether same phenomena occur for other loss functions in particular for cross-entropy. Second, this work focuses on unregularized regression tasks and it is important to identify optimal regularization schemes for pruning purposes. For instance, should we use classical $\ell_1/\ell_2$ regularization or can we refine them by injecting problem priors such as covariance information? Finally, going beyond pruning, using DC, one can investigate other compression techniques that process the output of the initial overparameterized learning problem, such as model quantization and distillation. 